\DocSection{Experimental Design}
\label{sec:experiments}

The empirical study separates four questions: (i) does the implementation
estimate the stated objective correctly, (ii) what approximation or optimization
error remains, (iii) how well does the learned distribution control the
environment under each execution rule, and (iv) what simulator and compute cost
is required? Direct variational learning is the practical candidate. Sampling
alternatives additionally require target-fidelity and mixing diagnostics;
their failures are not silently reinterpreted as successful inference.

\DocSubsection{Domains and scope}

Stochastic GridWorld is the planned primary study. Easy, medium, and hard
families use, respectively, \(3\!\times\!3\), \(5\!\times\!5\), and
\(8\!\times\!8\) grids. Wall counts are 1, 4, and 13; swamp counts are 1, 3,
and 9; gravel and dirt counts are each 1, 3, and 9. Relative to all grid cells,
wall fractions are 11.1\%, 16.0\%, and 20.3\%, and each of the other negative
terrain fractions is 11.1\%, 12.0\%, and 14.1\%. Action success probabilities
are 0.90, 0.85, and 0.80, with equal left/right slip; horizons are 20, 40, and 80.
Frozen-map audits record reachability, shortest-path length and count, and
single-cell start--goal bottlenecks. Map zero is used for development calibration;
maps one and two are reserved for fresh-map validation after configuration
selection. The confirmatory scientific design was frozen before generating
ten fresh base geometries per tier. Each base is paired with its meadow
representation and three fresh training blocks, yielding 30 independent base
geometries and 90 paired geometry/training blocks. Maps are accepted only for
the declared terrain counts, safe connectivity, and absence of prior or new
spatial duplicates, never according to any policy's performance. Small
diagnostic instances retain exact enumeration;
larger instances retain exact dynamic-programming evaluation of any fixed
policy even when enumerating all policies is impossible.

\begin{figure}[t]
  \centering
  \includegraphics[width=\linewidth]{figures/terrain-gridworld-tiers-20261003.pdf}
  \caption[Paired terrain layouts for the three GridWorld tiers]{Illustrative
  development layouts for the three GridWorld tiers. Navigation (top) and
  meadow reward landscapes (bottom) share base geometry, so they are not
  independent map replicates. S marks the start and G the absorbing goal.
  These layouts illustrate tasks, not experimental performance. Protected
  layouts remain unused during method development and tuning.}
  \label{fig:gridworld-tiers}
\end{figure}

Rewards accrue on arrival: pavement 0, gravel $-1$, dirt $-0.5$, meadow $+0.5$,
goal $+5$, and swamp $-5$, with an additional $-0.1$ per active transition.
The goal and swamp absorb; their padded suffix earns zero. A blocked move stays
in place and pays the actual destination terrain again. The paired reward
landscape adds 1, 2, and 3 nonterminal meadow cells without changing the base
geometry. Repeated meadow rewards are intentional: goal rate is not a utility
selection criterion in this secondary panel. All seven cell types are therefore
covered, without pooling navigation and meadow returns on an incompatible scale.

Only domains that add a distinct scientific stress test are considered after
GridWorld passes its medium gates:
\begin{table}[t]
  \centering
  \small
  \caption{Domain-extension decision.  Extensions are conditional on the
  primary GridWorld study rather than required to enlarge the result table.}
  \label{tab:domain-scope}
  \begin{tabular}{@{}p{25mm}p{62mm}p{35mm}@{}}
    \toprule
    Domain & Added evidence & Decision \\
    \midrule
    Blackjack & Compact nonspatial stochasticity and intrinsic return
      variance with two actions \citep{towersEtAl2025} & First low-cost
      extension after GridWorld \\
    Tireworld & Stochastic failures, resources, sparse success, and competing
      routes \citep{bonetGeffner2005} & Best second headline domain if time and
      compute remain \\
    Taxi & Action masks and subgoals, but mostly deterministic and spatial &
      Optional smoke test only \\
    Academic Advising & Factored multivalued state and concurrent actions
      \citep{guerinEtAl2012} & Future structured-policy work \\
    Continuous control & Requires a new reference measure and policy-space
      sampler & Out of scope \\
    \bottomrule
  \end{tabular}
\end{table}

A Tireworld extension would use a verified factored-state adapter; pyRDDLGym
is a possible RDDL-to-simulator route \citep{taitlerEtAl2022}.  Transition,
termination, reward, and action-mask tests on a small instance are mandatory
before any such results are included.

\DocSubsection{Competitors and controls}

The competitor implementations preserve each published core update while
using a common tabular code path for identical environment semantics, seeding,
transition accounting, and artifacts.  They are named by the exact implemented
variant rather than by a broader software package.

\begin{table}[t]
  \centering
  \footnotesize
  \caption{Learned competitors.  Uniform random behavior and exact
  finite-horizon dynamic programming are controls, not competitors.}
  \label{tab:competitors}
  \begin{tabular}{@{}p{26mm}p{62mm}p{37mm}@{}}
    \toprule
    Reported method & Implemented variant & Role \\
    \midrule
    REINFORCE & Monte Carlo policy gradient with learned state-value baseline
      \citep{williams1992} & On-policy score-function baseline \\
    PPO-Clip & Clipped surrogate, Monte Carlo advantages, entropy bonus; no GAE
      or value clipping \citep{schulman2017} & On-policy clipped-surrogate baseline and
      separately costed proposal trainer \\
    Discrete SAC & Categorical actor, twin critics, target critics, replay, and
      fixed entropy temperature \citep{christodoulou2019} & Off-policy
      entropy-regularized baseline \\
    Categorical CEM & Distribution search over complete stationary
      deterministic policies \citep{mannorEtAl2003} & Direct black-box policy
      search \\
    Double Q-learning & Online tabular Double Q-learning with linearly annealed
      epsilon-greedy exploration \citep{vanHasselt2010} & Value-learning
      baseline \\
    \bottomrule
  \end{tabular}
\end{table}

The exact control is nonstationary and may use the known model, so it is an
upper reference rather than a fair learned competitor.  All learned and
inferred output policies use stationary physical-state inputs.  For correct
finite-horizon backup targets, discrete SAC uses time-indexed twin critics and
replay tuples, while Double Q-learning uses time-indexed tables and projects its
result into the declared stationary policy class using state--time visitation.

Comparisons are stratified by information access.  In the known-model stratum,
exact-value tempering is tested for target fidelity and mixing against exact
enumeration or known-model controls. Fixed-tape scoring alone does not imply
simulator-only access: the current path/recovery block builder reads exact
successor probabilities. Such runs remain model-assisted even when scores
come from rollouts. In a simulator-only stratum, that builder must be disabled
or replaced by separately counted simulator-based construction. Eligible
fixed-tape variants, direct variational learning with its PPO initializer, and
the five learned competitors receive generative-simulator access. Exact-model
tempering may be displayed beside simulator-only learners as a descriptive
reference, but such a cross-stratum table cannot support a sample-efficiency
superiority claim.  Model value evaluations, simulator transitions, and guide
training transitions are therefore reported in separate cost columns.

\DocSubsection{Common evaluation semantics}

For the committed-policy diagnostic, the evaluator samples a complete
deterministic policy from the reported distribution and commits to it for the
entire episode.  For an explicit particle or MCMC sample, that complete policy
is used directly.  For statewise categorical outputs, one action is sampled
once at every decision state before evaluation. The fixed table's expected
return and goal probability are evaluated exactly; their distributional means
are estimated with 2,048 independent table draws at the terminal checkpoint.
The original 64-draw development learning checkpoints remain unchanged and
their first 64 terminal draws are preserved in the higher-precision attachment.
For independent VI or baseline policy draws, the sample standard
deviation describes policy variation and the standard error describes
conditional evaluation precision. Neither measures uncertainty across training
runs. Correlated MCMC draws or resampled particles require separate dependence
and ancestry diagnostics; the independent-draw formula is not valid for them.

This rule is an intervention on stochastic baseline outputs, not their native
deployment semantics. Practical comparisons and selection use method-native
expected original return: commitment for the inference distributions and CEM;
per-visit action draws for PPO, REINFORCE, SAC, and stochastic tie-breaking in
Double Q. Committed-policy values of stochastic baselines answer a narrower
policy-class question and cannot replace their native result.
A sampler's statewise marginals cannot replace its committed-policy result
because they remove policy-level dependence. The statewise
maximum-probability policy is also reported separately. Secondary metrics
include actual episode-return lower tails (not a histogram of policy expected
returns), separately counted evaluation interactions,
distributional target error where enumeration is possible, and mode recovery.

\DocSubsection{Fair budgets and tuning}

Every comparison uses the same maps, training seeds, evaluation streams,
horizon, reward, and termination convention.  Current GridWorld runs use
\(\gamma=1\); discounted baseline semantics must be validated before any
\(\gamma<1\) experiment.  Method-specific tuning grids are allowed because a
single global learning rate is not a fair tuning policy. Each tunable entry
receives its canonical setting and two prospective single-parameter alternatives
on all six panel/tier strata at 100k transitions. Actual allocated CPU costs are
disclosed and bounded, not assumed equal between learners. Exactly matching
canonical runs may be reused, with their original compute retained in the
ledger. Tuning maps and seeds
are disjoint from the final confirmatory set.

The inverse temperature defines the target and is not selected after observing
held-out control performance.  Exact diagnostic maps use a predeclared
sensitivity grid; the headline GridWorld target freezes one value on map-zero
development data. The prospective pure-inference screen contrasts \(\beta=64\)
with \(\beta=16\), or changes policy batch size from 32 to 8 without changing
temperature. Initialized inference contrasts guide fractions 0.5 and 0.75,
or the same batch-size change. One configuration per entry is selected over
equally weighted tier/panel normalized return gaps, not one winner per maze.
Every configuration gets the same 2,048-table terminal readout before selection.
This is a bounded target/configuration choice on development data, never a
best-of-grid choice using protected performance.

All interactions used by PPO guide training, variational refinement, SMC,
policy-bank evaluation, and sample-average MCMC count toward the proposed method.  Exact policy-value
evaluations, optimizer updates, wall time, CPU or accelerator time, and peak
memory are reported separately.  A generally applicable exploration aid is
offered to compatible methods or labeled as a separate assisted ablation.
Competitor bugs are fixed; no baseline is weakened to improve a comparison.

New variational and baseline runs enable \texttt{strict\_transition\_budget}:
the declared allowance is a maximum for all training transitions, including
the initializer. Refinement receives the total minus actual guide usage.
Complete-return reservation can leave a small unused allowance, which is
reported. Historical tempering configurations retain different semantics:
their \texttt{transition\_budget} covers guide training, and subsequent sampling
adds cost. They are not relabeled as strict-cap comparisons. Full-budget PPO
under the same allowance is an essential attribution control; the current
control starts from scratch, rather than resuming the guide optimizer.

\DocSubsection{Cost-effective execution order}

After the claim-specific checks in Section~\ref{sec:validation}, small
technical runs verify every learner's artifacts, budgets, and environment
semantics. A paired budget-response check informs resource requests without
dropping weak methods. The bounded 100k screen precedes 200k validation on
two fresh base geometries per tier, two independent training seeds per geometry,
and both terrain panels, retaining every declared competitor and control.
The two panels share geometry and seed blocks; they do not double independent
map replication. Configuration rankings may change between 100k and 200k.

Only after this validation are the inference engine, confirmatory size, cap,
analysis, and exclusions frozen before protected outcomes. Running jobs are
limited to two, at most one SAC job, using separate serial lanes. Independent
numeric and scheduler checks precede phase expansion. Technical failures halt
expansion; low returns do not justify excluding a method or layout. All attempts
are preserved, and a retry requires an explicit new artifact identity rather
than overwriting completed or partial results.

The frozen main engine is PPO-initialized variational inference; pure inference
and the full-budget assisted PPO control remain in the comparison. Every
learned run receives the same 200k total active-interaction ceiling and a
2,048-table terminal readout. A separate 4,096-episode native deployment
evaluation records actual returns, goal arrival, swamp termination, and timeout.
Random, return-optimal and goal-optimal controls are evaluated once per
map/panel, not duplicated for every training seed. Sampling and independent
scalar replay run on compute nodes; their actual interactions and allocated
CPU costs are kept separate from learning costs.

The primary report preserves original expected returns for each of the six
tier/panel strata, averaging training blocks within maps and maps equally.
The secondary composite gives equal weight to those six strata after
normalization between each map's random and return-oracle anchors; lower gap
is better. Degenerate anchor gaps retain raw results and undefined normalized
values, without clipping or redrawing. Joint hierarchical bootstrap resamples
base geometries, then paired training blocks, together across methods and
panels (10,000 draws, seed 0). Its 95\% intervals are descriptive and marginal,
not simultaneous tests or proof of equivalence. Episode-tail figures use the
predeclared medium/navigation and hard/meadow categories and aggregate all
their maps and training runs, rather than selecting attractive instances.

The same order applies to domain extension: one exact Blackjack instance, one
seed per method, then three-seed timing, and only then a complete table.
Tireworld work begins only after the GridWorld confirmatory manifest is frozen.

\DocSubsection{Current empirical status}

The redesigned terrain confirmatory study remains incomplete; no pilot score
is a final result. The numerical material below records HISTORICAL development
diagnostics from the earlier larger-grid study (8, 16, and 32 cells per side).
Those maps, budgets, and seed cohorts are different from the current 3/5/8
terrain design. They are retained for methodological context, not pooled into
its primary tables or used as protected evidence. Final results in the thesis
and compact paper will use the same complete independently verified data.

\DocSubsubsection{Variational allocation and full-budget controls}

A cached medium-map PPO guide trained for 200,015 transitions initialized three
independent refinements per rollout setting. With approximately one million
additional transitions, \(K=1\) produced committed returns 0.581, 0.595, and
0.548; \(K=4\) produced -0.144, -0.244, and -0.087. The corresponding update
counts were 662--669 versus 157--160. Every refinement improved its paired guide
estimate and estimated ELBO. Evaluation used 512 independent policy draws,
coupled within each guide/refinement pair. These refinements share one guide
and map, so they do not establish robustness across independently trained guides.

\input{shared/results/medium-development-table}
Fresh-guide, strict-cap runs then used nominally 200k guide transitions and 1M
refinement transitions. On the easy map, committed value rose from the guide's
3.010 to 3.318, below the nonstationary upper control 3.451. Actual usage was
1,199,880 transitions. The medium run used 1,199,725 transitions and achieved
committed value 0.725 versus its guide's -0.225.
Table~\ref{tab:medium-development} includes full-budget PPO controls at the same
medium training seed 572. The VI guide used count-bonus coefficient 1.0;
refinement used original rewards. Assisted PPO used coefficient 1.0 throughout,
and standard PPO used zero. Both controls trained from scratch.

VI's committed value is slightly above assisted PPO's and its native value is
slightly below. The committed difference 0.069 is smaller than the combined
conditional evaluation standard error. Neither equivalence nor superiority is
established. The shorter local runtime is likewise one observation, not a
general speedup claim. These runs share executable snapshot and map; their
wall times include offline exact evaluation and use one numerical-library thread.

\DocSubsubsection{Hard-map exploration and execution semantics}

\input{shared/results/development-comparison-figure}
Figure~\ref{fig:development-execution} displays the same completed development
runs under both execution rules. The hard-panel allocation change is separately
labeled; it is not an increase in the total simulator allowance.

The unchanged 200k-guide allocation failed on hard map zero with seed 573:
all guide and refinement batches returned the no-goal value -38.4. The estimated
ELBO rose through reduced KL, without improved control. The map is not
unreachable: its finite-horizon upper return control is -4.001.

A separately frozen two-step diagnostic first trained full-budget assisted PPO.
Its first successful batch occurred near 538k transitions, beyond the original
guide allowance. Native goal probability reached 0.995. This satisfied the
predeclared condition for testing a 1M-guide/200k-refinement reallocation at the
same total cap, map, and seed. That run consumed 999,721 guide and 199,549
refinement transitions. It achieved committed return -19.219 (conditional MCSE
0.679), compared with its guide's -23.077. Native return was -6.775.

Committed goal probability reached 0.700, versus zero with the short guide;
native action resampling reached 0.993. Full-budget assisted PPO achieved
committed return -20.900 and native return -6.415, with goal probabilities
0.637 and 0.995 respectively. Wall times were 132.4 seconds for reallocated
VI and 133.9 seconds for assisted PPO, including evaluation. This single-seed
allocation diagnostic supports testing independent guide seeds, not a general
hard-grid success or superiority claim. The large execution gap is also a
counterexample to treating policy commitment itself as a control advantage.

\DocSubsubsection{Independent-guide BGU calibration}

\input{shared/results/cluster-calibration-table}
\input{shared/results/cluster-calibration-figure}
The next frozen screen completed all 16 runs on BGU compute nodes. Training
seeds 601--603 were used on medium and hard development map zero, and seed
604 on separately reserved development map 100 in each tier. Held-out maps
1--3 remained unused. The total cap was 1.2M transitions per run, with nominal
VI guide allocations of 200k on medium and 1M on hard. The target, learning
rate, count bonus, and $N=32$, $K=1$ settings were fixed before dispatch.
PPO used count bonus 1 throughout; VI used it only in initialization and
unshaped rewards during refinement. These are explicitly assisted comparisons.

Every artifact passed checks of configuration, map hash, executable snapshot,
strict budget, finite oracle-bounded metrics, and Slurm provenance. The cluster
cohort shares its software environment and is not pooled with older local
timing or seed results. Evaluation uses 256 independent complete-policy draws;
their conditional Monte Carlo precision is not uncertainty across training runs.
The plotted arithmetic means in Figure~\ref{fig:cluster-calibration} retain
every seed, including the failed medium run. Table~\ref{tab:cluster-calibration}
also reports the second-map checks separately rather than blending maps.

PPO has higher committed and native return in seven of the eight paired
map--seed comparisons. Primary-map mean committed returns are -4.973 versus
0.449 on medium and -17.597 versus -15.493 on hard (VI versus PPO).
The corresponding native means are -4.695 versus 0.970 and -6.884 versus
-6.076. The second development maps also favor PPO in return under both rules.
This does not establish a universal ranking, but it contradicts extrapolating
the earlier favorable single-seed pilot into a robustness claim.

Medium seed 602 isolates the failure. Its guide stops at 199,923 transitions
with native return -15.936 and committed return -15.999. Every one of its 196
refinement minibatches returns -16 with essentially zero between-policy return
variation. The estimated ELBO rises from -1140.570 to -1026.622 through the
entropy term, not successful navigation. Final native goal probability is
0.00371; committed goal probability is approximately $3.1\times10^{-6}$.
Continued assisted PPO with the same seed reaches its first positive batch
mean at 318,132 transitions. Some guide batches had successful trajectories
earlier, so first-ever success is not by itself a sufficient handoff criterion.

The primary medium-map mean wall times are 40.5 seconds for VI and 130.4 for
PPO; hard-map means are 170.0 and 174.6 seconds. They include offline evaluation
and describe this CPU implementation, not hardware-independent algorithmic
complexity. A shorter failed run is not a performance advantage. A separate
hard-grid SAC throughput screen consumed 4,992 transitions in 325.56 seconds;
its larger conditional profile was therefore not released. That tiny-budget
profile is cost evidence only, not a fair SAC control-performance comparison.

\DocSubsubsection{Bounded allocation follow-up}

\input{shared/results/allocation-followup-table}
\input{shared/results/allocation-followup-figure}
A separately frozen follow-up reallocated the medium allowance to 600k guide
and the remainder refinement, without changing either algorithm's updates or
increasing total interactions. Selected failure seed 602 recovered to committed
goal probability 0.9863 and native probability 0.9996. This satisfied the
predeclared conditions for two fresh paired seeds, 606 and 607, and a second-map
VI check using seed 604. The repair is not pooled into the fresh-seed estimate.

Both fresh VI runs avoid collapse, with committed goal probabilities 0.9910
and 0.9779; native probabilities exceed 0.9998 for both methods. Committed
returns are 0.628 and -0.004 for VI versus 0.626 and 0.258 for PPO. Native
returns are 0.975 and 0.737 versus 1.055 and 0.874. Thus the improvement removes
the observed exploration failure but does not establish superiority: PPO still
has higher native return on both fresh seeds. On development map 100, VI has
committed/native returns 0.752/1.050 versus the existing PPO control's
0.781/1.048. These near values are descriptive, not an equivalence test.

All six completed follow-up artifacts passed the integrity audit. Two PPO
attempts were interrupted by an unhandled Slurm warning signal; the launcher
was corrected and only those attempts were repeated from scratch at the same
five-minute wall limit. Failed logs and partial artifacts are retained and
their wasted compute is counted separately. No completed scientific outcome
was replaced. Follow-up timings come from a different node than the original
calibration; only the within-follow-up fresh pairs have matched hardware.
Their mean wall times, including evaluation, are 102.5 seconds for VI and
182.0 for PPO. This small CPU comparison motivates a costed larger study,
not a general speedup claim. Held-out maps remain unused.

\DocSubsubsection{Sampling alternatives: development findings}

Preliminary medium-GridWorld exact-value
tempering diagnostics produced high-return committed policies but exposed a sharp
intermediate-temperature exchange bottleneck.  Four independent 32-temperature,
power-four chains, including an overdispersed chain, each retained 800 draws.
They failed the predeclared gate: split return \(\widehat R=1.096\), minimum
per-chain return ESS of 10.5, minimum whole-run adjacent-swap acceptance of 0.049, no
completed round trip, and maximum pairwise state-action total variation of 0.699.
The principal warm-chain bottleneck lay between inverse temperatures 1.44 and
1.98.  Because the legacy acceptance counter included burn-in, subsequent runs
also store a separate post-burn rate; the ESS, round-trip, and between-chain
failures do not depend on that counter.  A subsequent two-chain,
4,000-iteration custom-ladder probe also failed.  Its warm chain had
post-burn acceptance between 0.598 and 0.997 on every edge, but no walker
traversed both endpoints.  In the overdispersed chain, replicas below the cold
endpoint remained on a flat random-policy plateau near return \(-16\), while
the cold replica remained near return \(1.1\); the final edge consequently
accepted no swap.  High acceptance inside the bad plateau therefore did not
indicate useful mixing.

Subsequent whole-policy and structural-block rescue attempts did not establish
adequate mixing. An isolated path-proposal screen used four independent
4,500-step chains at \(\beta=3.42\), discarding 1,200 steps and retaining
every tenth subsequent state. Low-initialized chains remained near return
\(-16\); high-initialized chain medians were approximately \(-1.56\) and
\(-1.41\). Path banks fitted on one chain from each region failed to generate
useful bidirectional crossings on the other chains. These measurements do not
estimate the relative posterior mass of the regions.

A bounded three-seed guide-field replica-exchange screen subsequently found
high-return policies but did not validate sampling adequacy. Return effective
sample sizes were 28--66, split \(\widehat R\) was 1.032, maximum between-chain
coordinate total variation was 0.578, and completed round trips were 0, 1, and 1.
A fixed-bank SMC screen retained only 3, 4, and 4 root ancestors from 256
particles and returned values near -16. These alternatives have different
costs and information access; they are not pooled into a simulator-efficiency
leaderboard. Historical failed checks remain failed. The practical VI candidate
instead requires family-error, learning-stability, and execution-aware control
checks before full-scale confirmation.

\begin{table}[t]
  \centering
  \caption{Required confirmatory reporting.  No entry is populated from pilot
  runs.}
  \label{tab:result-matrix}
  \begin{tabular}{@{}p{25mm}p{42mm}p{55mm}@{}}
    \toprule
    Question & Comparison & Primary report \\
    \midrule
    Estimator correctness & exact differentiation & gradient error under stochastic visitation \\
    Family error & enumerable posterior & variational gap and omitted dependence \\
    Sampling alternatives & chains and ladders & ESS, swaps, ancestry, and agreement \\
    Control quality & five competitors & native and committed return and goal probability \\
    Robustness & map families and seeds & uncertainty and failed-run count \\
    Scalability & matched accounting & interactions, time, and memory \\
    \bottomrule
  \end{tabular}
\end{table}
